\chapter{Discussion}
\label{ch:discussion}

The evaluation separated results that hold regardless of how activations were simulated from
results that depend on it. This chapter examines what follows for deployment and for the next
phase. \Cref{sec:disc:validity} discusses threats to validity. \Cref{sec:disc:adaptive}
considers an adaptive attacker. \Cref{sec:disc:toctou,sec:disc:amortise} address the timing and
cost of attestation. \Cref{sec:disc:llm} discusses what changes for real language models, and
\cref{sec:disc:labels} the role of labels. \Cref{sec:disc:ethics} closes with ethical
considerations.

% ---------------------------------------------------------------------------
\section{Threats to validity}
\label{sec:disc:validity}

\paragraph{Construct validity: simulated activations.} The largest threat is that activation-level
attacks are defined as Gaussian distributions that differ from the clean Gaussian in mean or
spread. The monitor models the clean Gaussian. Detection of such attacks is close to guaranteed
whenever the difference is large, and impossible when it is small, so the measured rates reflect
choices made in the simulator. Real activations are not Gaussian. Transformer activations in
particular are known to contain large outliers in a few channels, the problem SmoothQuant
addresses \cite{xiao2023smoothquant}. Such outliers would inflate the clean $z$ distribution and
force higher thresholds, and per-channel normalisation would be essential.

\paragraph{Internal validity: shared authorship.} The detector, the simulators and the evaluation
were written together. The two defects of \cref{sec:impl:bugs} that sat in the simulators, and
the covariance defect that sat in the detector, show that such a pipeline can hide errors.
Independent references in the tests mitigate this risk, and the fact that the clean false-positive
rate matches a closed-form prediction is reassuring, because a gross error in the scoring path
would break that agreement.

\paragraph{External validity: stub weights and the test machine.} The weight sets have
realistic tensor shapes but random contents. Attestation cost depends only on byte count, so the
scaling results transfer to real checkpoints of the same size. They do not transfer to the NPU
or to a phone CPU, whose SHA-256 throughput may be higher, with hardware extensions, or lower.

\paragraph{Statistical conclusion validity.} With \NPerScenario{} inputs per scenario, a measured
rate of 100\% has a one-sided 95\% lower confidence bound of about 99.7\%, by the ``rule of
three''. A measured rate of 0\% has an upper bound of about 0.3\%. These bounds are adequate for
the claims made. The false-positive comparison (predicted \PredFPRZFour{}, measured \EmpFPRZFour{}
with a seed-to-seed standard deviation of 1.5 points) is consistent within sampling error.

% ---------------------------------------------------------------------------
\section{An adaptive attacker}
\label{sec:disc:adaptive}

Suppose the attacker knows the calibration statistics and thresholds. The two detectors respond
very differently.

\paragraph{Against attestation,} knowledge does not help. Evading it requires a SHA-256 second
preimage for a modified tensor, a forged ECDSA signature, or control of the verifier, and all
three are excluded by the threat model. The realistic adaptive strategies are temporal rather than
cryptographic: flip bits just after a verification and restore them before the next one, or
attack between verifications (\cref{sec:disc:toctou}).

\paragraph{Against the monitor,} knowledge helps a great deal. An attacker optimising an
adversarial input can add a penalty that keeps every monitored channel's $|z|$ below $\tau_z$ and
every Mahalanobis distance below $\tau_D$. Whether a successful attack remains feasible under that
constraint depends on how much freedom the input space leaves once activations are confined to the
calibration region. Carlini and Wagner showed that ten published adversarial-example detectors could
each be bypassed by attacks constructed with knowledge of the detector
\cite{carlini2017bypassing}. The honest default assumption is that a single-statistic monitor will not stop a determined adaptive
attacker. Its realistic roles are to raise the cost of attacks, to catch non-adaptive ones such as
off-the-shelf jailbreak prompts and faults, and to provide forensic signal.

Three directions could make the monitor harder to evade, and each is a candidate for the next
phase. First, keeping the set of monitored layers or channels secret and rotating it would force
the attacker to satisfy constraints on everything. Second, statistics that are expensive to
satisfy jointly, such as the joint distribution across layers, could replace per-layer marginals.
Third, the monitor could be coupled to the model's output, so that inputs that change the output a
great deal while keeping activations typical become suspicious.

% ---------------------------------------------------------------------------
\section{Time of check versus time of use}
\label{sec:disc:toctou}

Attestation certifies the weights at the moment they are hashed. If verification runs at times
$t_1<t_2<\dots$ and inference reads weights at $t$, a flip at $t_1<t'<t$ affects that inference
undetected, and it is caught only at the next verification. Two quantities therefore matter: the
\emph{detection latency}, the time from a flip to the next verification that covers the modified
tensor, and the \emph{exposure}, the number of inferences that run on modified weights in the
meantime.

Closing the gap completely would require verifying the bytes on the path into the accelerator,
for example by hashing in the DMA engine or in on-chip SRAM after transfer. That is a hardware
change. For a software-only design, the realistic goal is a bounded
detection latency at an acceptable cost, the subject of the next section.

% ---------------------------------------------------------------------------
\section{Amortised attestation}
\label{sec:disc:amortise}

\Cref{sec:eval:merkle} measured \MerkVerLlamaPB{}\,ms to verify the 1.4\,GiB Llama-3.2-1B stub.
Suppose a budget of $b$ seconds of hashing per inference step, and hashing throughput $\rho$
bytes per second. Each step can verify $b\rho$ bytes. The Merkle tree makes partial verification
sound: any subset of tensors can be checked against the signed root using their membership
proofs (\cref{fig:merkle}), without hashing the rest.

\paragraph{Round-robin verification.} Verify tensors in a fixed cyclic order, $b\rho$ bytes per
step. A full cycle covering all $B$ bytes takes $\lceil B/(b\rho)\rceil$ steps, and every
modification is detected within one cycle. At the measured $\rho\approx\MerkThrLlamaPB$\,MiB/s and a budget of
$b=5$\,ms per token, $b\rho\approx 4.5$\,MiB per token, and the Llama-3.2-1B stub
($B\approx\MerkSizeLlamaPB$\,MiB) is fully covered every $\approx315$ tokens. An attacker who
knows the schedule can, however, time a flip-and-restore to avoid it.

\paragraph{Randomised verification.} Choose each step's tensors at random, weighted by size, with
the choice secret. A modified tensor of $\beta$ bytes is covered in a given step with probability
about $q = \min(1, b\rho/B)$ if verification is byte-proportional, and it remains undetected
after $t$ steps with probability $(1-q)^t$. With the numbers above, $q\approx 1/315$. A
modification persisting for 1{,}000 tokens escapes detection with probability
$(1-1/315)^{1000}\approx 4\%$, and one persisting for 2{,}000 tokens with probability about
0.2\%. Randomisation removes the attacker's ability to predict the schedule at the cost of a
probabilistic rather than deterministic bound. Combining the two, a random permutation per cycle,
gives a deterministic bound of one cycle with unpredictable timing within it.

\paragraph{Prioritised verification.} Not all tensors matter equally. Hong et al.\ found that a
large fraction of single-bit vulnerabilities concentrate in particular parameters
\cite{hong2019terminal}, HASHTAG selects layers by measured sensitivity
\cite{javaheripi2021hashtag}, and NeuroPots steers attacks into known trapdoors
\cite{liu2023neuropots}. Verifying sensitive tensors more often shortens their detection latency
at no extra cost. For transformers, the embedding and output projections are a large share of
the bytes ($2VH$ in \cref{eq:stub_params}, about 36\% of the Llama-3.2-1B stub's parameters). Their
sensitivity to single flips relative to the attention and MLP projections is an open empirical
question that the next phase can measure on a real model.

\paragraph{Faster hashing.} SHA-256 throughput on CPUs with SHA extensions is typically several
times higher than the \MerkThrMin{}--\MerkThrMax{}\,MiB/s measured on this container. A non-cryptographic
checksum as a first-line filter, with cryptographic verification only on a schedule, is the
trade-off RADAR and HASHTAG make \cite{li2021radar,javaheripi2021hashtag}. Such a filter must be
keyed, or its positions secret, because an attacker who knows an unkeyed checksum can flip
compensating bits. A keyed MAC computed with a key held in the TEE retains cryptographic strength
and is typically faster than a hash plus a signature check.

% ---------------------------------------------------------------------------
\section{From stubs to real language models}
\label{sec:disc:llm}

Moving the monitor to a real quantized LLM changes four things.

\begin{enumerate}[leftmargin=*]
  \item \textbf{Channel count.} A model with hidden size 2048 and 16 blocks, monitored at two
        points per block, has $m\approx 65{,}000$ channels per token. \Cref{eq:bonf} gives
        $z^*\approx 5.25$ for $\alpha=1\%$, compared with \ZBonf{} here. The correction remains
        cheap to compute, but every attack signal must exceed a higher bar.
  \item \textbf{Covariance.} A dense $2048\times2048$ covariance per monitoring point is 32\,MiB
        in double precision. Diagonal statistics, low-rank-plus-diagonal models, or projection
        onto a few principal components are the realistic options.
  \item \textbf{Tokens.} A language model produces one activation vector per token. Monitoring
        can score each token, which is expensive and multiplies the comparisons, or summarise a
        sequence. Prompt-level attacks are naturally scored per sequence, generation-time
        extraction per token.
  \item \textbf{Validating the memorisation hypothesis.} The extraction-probe detector assumes
        that regurgitating memorised text produces confident, low-variance late activations. A
        direct test is possible with public tools. Take a small open model, a set of known
        memorised sequences such as those identified by extraction studies, and matched
        non-memorised controls, and compare late-layer activation statistics while the model
        generates them. If no difference appears, the memorisation score should be dropped. If
        one does, its thresholds can be fitted to data rather than to the simulator.
\end{enumerate}

% ---------------------------------------------------------------------------
\section{On threat labels}
\label{sec:disc:labels}

The classifier's labels were useful in the evaluation mainly because they failed in informative
ways. Steganographic poisoning and runtime bit flips cannot be told apart from the weights alone.
Inputs that disturb every layer were called faults. And an entropy score invariant to shifts and
scaling could not respond to the simulated jailbreak. Each failure pointed to missing evidence
rather than to a tuning problem. Provenance requires evidence over time, the fault-versus-attack
distinction requires a comparison that is not biased by layer counts, and jailbreak detection
requires a score sensitive to the property of interest.

A more robust design would make the label a function of explicitly listed evidence, such as
``hash mismatch at load time'', ``hash mismatch first seen at runtime'', ``anomaly rate in early
layers $>$ late layers'' and ``memorisation score above threshold''. The output would report the
evidence alongside a suggested interpretation, rather than collapsing everything into one of eight
names. The explanation strings already list most of this evidence. Promoting it to structured
output would let operators write their own response policies.

% ---------------------------------------------------------------------------
\section{Ethical considerations}
\label{sec:disc:ethics}

The attack simulators in this project do not implement working attacks. The bit-flip simulators
modify arrays in memory and contain no Rowhammer code. The extraction and probe simulators
generate synthetic activation distributions and recover no data. Nothing in the repository
lowers the barrier to attacking a real system beyond what the cited literature already
describes. The defensive components are intended to be published openly so that others can
evaluate and improve them.

Two downstream concerns deserve mention. A monitor that inspects activations sees, indirectly,
information about user inputs, so a deployment should not log activations or explanations
containing them without the same care applied to the inputs themselves. And false positives have
costs: a 7\% false-alarm rate that rejects legitimate requests would harm users. This is one
reason the threshold analysis of \cref{sec:eval:fpr} matters beyond the evaluation.
